\documentclass{article}
\usepackage[margin=20mm]{geometry}
\usepackage{surface-diagrams}
\pagestyle{empty}
\begin{document}
\begin{center}
\Large Thesis blue and gray\par\bigskip
\resizebox{0.95\linewidth}{!}{\SurfaceDiagram{planar-default.tikz}}
\end{center}
\clearpage
\begin{center}
\Large More room around the same five objects\par\bigskip
\resizebox{0.95\linewidth}{!}{\SurfaceDiagram{planar-large-ellipse.tikz}}
\end{center}
\clearpage
\begin{center}
\Large Twelve marked points\par\bigskip
\resizebox{0.95\linewidth}{!}{\SurfaceDiagram{planar-twelve-points.tikz}}
\end{center}
\clearpage
\begin{center}
\Large Twenty points; smaller dots\par\bigskip
\resizebox{0.95\linewidth}{!}{\SurfaceDiagram{planar-twenty-points.tikz}}
\end{center}
\clearpage
\begin{center}
\Large Six boundaries interspersed with points\par\bigskip
\resizebox{0.95\linewidth}{!}{\SurfaceDiagram{planar-mixed-order.tikz}}
\end{center}
\clearpage
\begin{center}
\Large Outer ellipse hidden\par\bigskip
\resizebox{0.95\linewidth}{!}{\SurfaceDiagram{planar-no-outline.tikz}}
\end{center}
\clearpage
\begin{center}
\Large Boundary radius 8; point radius 3\par\bigskip
\resizebox{0.95\linewidth}{!}{\SurfaceDiagram{planar-large-boundaries.tikz}}
\end{center}
\clearpage
\begin{center}
\Large Point radius 8; boundary radius 3\par\bigskip
\resizebox{0.95\linewidth}{!}{\SurfaceDiagram{planar-large-points.tikz}}
\end{center}
\clearpage
\begin{center}
\Large Configurable colors, white background\par\bigskip
\resizebox{0.95\linewidth}{!}{\SurfaceDiagram{planar-custom-colors.tikz}}
\end{center}
\clearpage
\begin{center}
\Large Free positions off the axis (no cut curves)\par\bigskip
\resizebox{0.95\linewidth}{!}{\SurfaceDiagram{planar-custom-positions.tikz}}
\end{center}
\clearpage
\begin{center}
\Large Individual radii 3, 5, 7, 9\par\bigskip
\resizebox{0.95\linewidth}{!}{\SurfaceDiagram{planar-per-object-size.tikz}}
\end{center}
\clearpage
\begin{center}
\Large An empty, rounder disk\par\bigskip
\resizebox{0.95\linewidth}{!}{\SurfaceDiagram{planar-empty.tikz}}
\end{center}
\clearpage
\begin{center}
\Large An empty disk in circular-boundary mode\par\bigskip
\resizebox{0.95\linewidth}{!}{\SurfaceDiagram{circles-empty.tikz}}
\end{center}
\clearpage
\begin{center}
An annulus with an arc joining its boundaries\par
SVG example available; circular-boundary TikZ support is planned for P7.
\end{center}
\clearpage
\begin{center}
Circular holes: default row and sizes\par
SVG example available; circular-boundary TikZ support is planned for P7.
\end{center}
\clearpage
\begin{center}
Circular holes and marked points\par
SVG example available; circular-boundary TikZ support is planned for P7.
\end{center}
\clearpage
\begin{center}
\Large The same boundaries as gray dots\par\bigskip
\resizebox{0.95\linewidth}{!}{\SurfaceDiagram{circles-dot-comparison.tikz}}
\end{center}
\clearpage
\begin{center}
A straight arc joins two circle rims\par
SVG example available; circular-boundary TikZ support is planned for P7.
\end{center}
\clearpage
\begin{center}
An arc above the row ends on the rims\par
SVG example available; circular-boundary TikZ support is planned for P7.
\end{center}
\clearpage
\begin{center}
A reversed arc below the row\par
SVG example available; circular-boundary TikZ support is planned for P7.
\end{center}
\clearpage
\begin{center}
An arc from a hole to a marked point\par
SVG example available; circular-boundary TikZ support is planned for P7.
\end{center}
\clearpage
\begin{center}
An arc from the outer boundary to a hole\par
SVG example available; circular-boundary TikZ support is planned for P7.
\end{center}
\clearpage
\begin{center}
A closed curve surrounds circular holes\par
SVG example available; circular-boundary TikZ support is planned for P7.
\end{center}
\clearpage
\begin{center}
Numbered horizontal guides with a cut-crossing arc\par
SVG example available; circular-boundary TikZ support is planned for P7.
\end{center}
\clearpage
\begin{center}
Individual circle radii: 8, 14, 20\par
SVG example available; circular-boundary TikZ support is planned for P7.
\end{center}
\clearpage
\begin{center}
Transparent holes on a colored background\par
SVG example available; circular-boundary TikZ support is planned for P7.
\end{center}
\clearpage
\begin{center}
\Large Arc(1, 6)\par\bigskip
\resizebox{0.95\linewidth}{!}{\SurfaceDiagram{curves-simple-up.tikz}}
\end{center}
\clearpage
\begin{center}
\Large Arc(1, 6, start\_up=False)\par\bigskip
\resizebox{0.95\linewidth}{!}{\SurfaceDiagram{curves-simple-down.tikz}}
\end{center}
\clearpage
\begin{center}
\Large Arc(2, 5, cuts=(3,), start\_up=False)\par\bigskip
\resizebox{0.95\linewidth}{!}{\SurfaceDiagram{curves-one-cut.tikz}}
\end{center}
\clearpage
\begin{center}
\Large Arc(2, 3, cuts=(4, 2, 4, 2, 4), start\_up=False)\par\bigskip
\resizebox{0.95\linewidth}{!}{\SurfaceDiagram{curves-repeated-cuts.tikz}}
\end{center}
\clearpage
\begin{center}
\Large Arc(3, 6, cuts=(4, 0, 4))\par\bigskip
\resizebox{0.95\linewidth}{!}{\SurfaceDiagram{curves-winding.tikz}}
\end{center}
\clearpage
\begin{center}
\Large Arc(0, 7): endpoints on the outer boundary\par\bigskip
\resizebox{0.95\linewidth}{!}{\SurfaceDiagram{curves-outer-endpoints.tikz}}
\end{center}
\clearpage
\begin{center}
\Large Arc(0, 4, cuts=(2,))\par\bigskip
\resizebox{0.95\linewidth}{!}{\SurfaceDiagram{curves-outer-to-point.tikz}}
\end{center}
\clearpage
\begin{center}
\Large Loop((1, 4)): around objects 2, 3, 4\par\bigskip
\resizebox{0.95\linewidth}{!}{\SurfaceDiagram{curves-consecutive-loop.tikz}}
\end{center}
\clearpage
\begin{center}
\Large Loop((0, 6, 5, 1)): around the two ends\par\bigskip
\resizebox{0.95\linewidth}{!}{\SurfaceDiagram{curves-nonconsecutive-loop.tikz}}
\end{center}
\clearpage
\begin{center}
\Large Nested loops plus an interior arc\par\bigskip
\resizebox{0.95\linewidth}{!}{\SurfaceDiagram{curves-nested.tikz}}
\end{center}
\clearpage
\begin{center}
\Large Three disjoint arcs\par\bigskip
\resizebox{0.95\linewidth}{!}{\SurfaceDiagram{curves-disjoint.tikz}}
\end{center}
\clearpage
\begin{center}
\Large Loop around boundaries and marked points\par\bigskip
\resizebox{0.95\linewidth}{!}{\SurfaceDiagram{curves-mixed-loop.tikz}}
\end{center}
\clearpage
\begin{center}
\Large Arc(3, 4): default is straight\par\bigskip
\resizebox{0.95\linewidth}{!}{\SurfaceDiagram{directions-adjacent-default.tikz}}
\end{center}
\clearpage
\begin{center}
\Large Arc(3, 4, direction="up")\par\bigskip
\resizebox{0.95\linewidth}{!}{\SurfaceDiagram{directions-adjacent-up.tikz}}
\end{center}
\clearpage
\begin{center}
\Large Arc(3, 4, direction="down")\par\bigskip
\resizebox{0.95\linewidth}{!}{\SurfaceDiagram{directions-adjacent-down.tikz}}
\end{center}
\clearpage
\begin{center}
\Large Arc(2, 5): default still goes up\par\bigskip
\resizebox{0.95\linewidth}{!}{\SurfaceDiagram{directions-nonadjacent-default.tikz}}
\end{center}
\clearpage
\begin{center}
\Large Arc(0, 1): straight from outer boundary\par\bigskip
\resizebox{0.95\linewidth}{!}{\SurfaceDiagram{directions-outer-default.tikz}}
\end{center}
\clearpage
\begin{center}
\Large Arc(0, 1, direction="up")\par\bigskip
\resizebox{0.95\linewidth}{!}{\SurfaceDiagram{directions-outer-up.tikz}}
\end{center}
\clearpage
\begin{center}
\Large Arc(0, 1, direction="down")\par\bigskip
\resizebox{0.95\linewidth}{!}{\SurfaceDiagram{directions-outer-down.tikz}}
\end{center}
\clearpage
\begin{center}
\Large Arc(6, 7): straight to outer boundary\par\bigskip
\resizebox{0.95\linewidth}{!}{\SurfaceDiagram{directions-right-default.tikz}}
\end{center}
\clearpage
\begin{center}
\Large Arc(2, 3, cuts=(4,)): itinerary preserved\par\bigskip
\resizebox{0.95\linewidth}{!}{\SurfaceDiagram{directions-adjacent-with-cuts.tikz}}
\end{center}
\clearpage
\begin{center}
\Large Straight chain with a curved arc over it\par\bigskip
\resizebox{0.95\linewidth}{!}{\SurfaceDiagram{directions-shared-endpoints.tikz}}
\end{center}
\clearpage
\begin{center}
\Large Consecutive boundary dots also join straight\par\bigskip
\resizebox{0.95\linewidth}{!}{\SurfaceDiagram{directions-boundary-dots.tikz}}
\end{center}
\clearpage
\begin{center}
\Large Default, up, down on three adjacent pairs\par\bigskip
\resizebox{0.95\linewidth}{!}{\SurfaceDiagram{directions-three-directions.tikz}}
\end{center}
\clearpage
\begin{center}
\Large Genus 1, no boundary components\par\bigskip
\resizebox{0.95\linewidth}{!}{\SurfaceDiagram{genus-genus-one.tikz}}
\end{center}
\clearpage
\begin{center}
\Large Genus 2, no boundary components\par\bigskip
\resizebox{0.95\linewidth}{!}{\SurfaceDiagram{genus-genus-two.tikz}}
\end{center}
\clearpage
\begin{center}
\Large Genus 3, with involution axis\par\bigskip
\resizebox{0.95\linewidth}{!}{\SurfaceDiagram{genus-genus-three.tikz}}
\end{center}
\clearpage
\begin{center}
\Large Genus 5\par\bigskip
\resizebox{0.95\linewidth}{!}{\SurfaceDiagram{genus-genus-five.tikz}}
\end{center}
\clearpage
\begin{center}
\Large Type I: slots 1 and 6\par\bigskip
\resizebox{0.95\linewidth}{!}{\SurfaceDiagram{genus-fixed-ends.tikz}}
\end{center}
\clearpage
\begin{center}
\Large Type I: slots 2 and 3\par\bigskip
\resizebox{0.95\linewidth}{!}{\SurfaceDiagram{genus-fixed-handle.tikz}}
\end{center}
\clearpage
\begin{center}
\Large All six Type I slots on genus 2\par\bigskip
\resizebox{0.95\linewidth}{!}{\SurfaceDiagram{genus-all-fixed.tikz}}
\end{center}
\clearpage
\begin{center}
\Large Two Type II pairs: four boundaries\par\bigskip
\resizebox{0.95\linewidth}{!}{\SurfaceDiagram{genus-paired-top-bottom.tikz}}
\end{center}
\clearpage
\begin{center}
\Large Left and right Type II pairs\par\bigskip
\resizebox{0.95\linewidth}{!}{\SurfaceDiagram{genus-paired-left-right.tikz}}
\end{center}
\clearpage
\begin{center}
\Large One Type I plus two Type II pairs\par\bigskip
\resizebox{0.95\linewidth}{!}{\SurfaceDiagram{genus-mixed-types.tikz}}
\end{center}
\clearpage
\begin{center}
\Large Six pairs on a wider genus-3 surface\par\bigskip
\resizebox{0.95\linewidth}{!}{\SurfaceDiagram{genus-many-pairs.tikz}}
\end{center}
\clearpage
\begin{center}
\Large D3/D2A: six top-bottom pairs and an end boundary\par\bigskip
\resizebox{0.95\linewidth}{!}{\SurfaceDiagram{genus-d3-default-boundaries.tikz}}
\end{center}
\clearpage
\begin{center}
\Large E2: side pairs and two top-bottom pairs\par\bigskip
\resizebox{0.95\linewidth}{!}{\SurfaceDiagram{genus-e2-default-boundaries.tikz}}
\end{center}
\clearpage
\begin{center}
\Large Four pairs: automatic positions and sizes\par\bigskip
\resizebox{0.95\linewidth}{!}{\SurfaceDiagram{genus-automatic-pairs.tikz}}
\end{center}
\clearpage
\begin{center}
\Large Optional wider spacing and taller body\par\bigskip
\resizebox{0.95\linewidth}{!}{\SurfaceDiagram{genus-custom-proportions.tikz}}
\end{center}
\clearpage
\begin{center}
\Large Transparent collars on a colored background\par\bigskip
\resizebox{0.95\linewidth}{!}{\SurfaceDiagram{genus-colored-background.tikz}}
\end{center}
\clearpage
\begin{center}
\Large View from below / right\par\bigskip
\resizebox{0.95\linewidth}{!}{\SurfaceDiagram{views-below-right.tikz}}
\end{center}
\clearpage
\begin{center}
\Large View from below / left\par\bigskip
\resizebox{0.95\linewidth}{!}{\SurfaceDiagram{views-below-left.tikz}}
\end{center}
\clearpage
\begin{center}
\Large View from above / right\par\bigskip
\resizebox{0.95\linewidth}{!}{\SurfaceDiagram{views-above-right.tikz}}
\end{center}
\clearpage
\begin{center}
\Large View from above / left\par\bigskip
\resizebox{0.95\linewidth}{!}{\SurfaceDiagram{views-above-left.tikz}}
\end{center}
\clearpage
\begin{center}
Two figures in one document\par\bigskip
\SurfaceDiagram[0.6]{directions-adjacent-up.tikz}\par
\SurfaceDiagram[0.6]{directions-adjacent-down.tikz}
\end{center}
\end{document}
